\chapter{Non-pluripolar products}
\label{chap:q09}

Let $X$ be a complex manifold of dimension $n\geq1$. The bounded Bedford--Taylor product can concentrate near the poles of an unbounded potential. Truncating the potentials and retaining only the common finite levels removes that concentration. The resulting measures increase; a local mass bound makes their limit a current. Closedness requires a separate estimate for the derivatives of the level cutoffs. On a compact K\"ahler manifold, bounded global potentials supply the mass bound through Stokes' formula. Sources and precise comparisons are in \cref{sec:history-nonpluripolar}.

\section{Truncation and plurifine locality}
\label{sec:nonpluripolar-construction}
On a coordinate domain $U\subset X$, let $u_1,\ldots,u_p$ be plurisubharmonic, $1\leq p\leq n$. Put
\[
 u_{j,k}=\max\{u_j,-k\},\qquad
 O_k=\bigcap_{j=1}^p\{u_j>-k\},\qquad
 S_k=\mathbf1_{O_k}\ddc u_{1,k}\wedge\cdots\wedge\ddc u_{p,k}.
\]
The bounded products are those of \cref{def:bedford-taylor-product}, and the indicators restrict their coefficient measures. The sets $O_k$ are Borel and plurifine open. They need not be open in the ordinary topology.

\begin{lemma}\label{lem:nonpluripolar-truncation-limit}
The positive currents with measure coefficients $S_k$ increase in the order of positive currents. They have a locally finite limit precisely when
\begin{equation}\label{eq:nonpluripolar-local-mass-condition}
 \sup_k\int_K S_k\wedge\beta^{n-p}<\infty
 \quad\text{for every }K\Subset U,
\end{equation}
where $\beta$ is the Euclidean K\"ahler form. In that case the limit $S$ has no mass on any pluripolar set, and $S_k\to S$ in coefficient variation on compact subsets. If Borel sets $E_k\subset O_k$ satisfy $\mathbf1_{E_k}\to1$ almost everywhere for the trace measure of $S$, then $\mathbf1_{E_k}S_k\to S$ against bounded Borel tests with compact support. In particular this holds for an increasing exhaustion whose complement is pluripolar.
\end{lemma}
\begin{proof}
For $l\geq k$, all corresponding bounded potentials agree on $O_k$. \Cref{thm:bt-plurifine-locality} gives $\mathbf1_{O_k}S_l=S_k$. On the complement of $O_k$, $S_k$ is zero and $S_l$ is positive. Thus $S_l-S_k\geq0$.

Their trace measures increase, and \eqref{eq:nonpluripolar-local-mass-condition} gives a locally finite measure limit. The positive coefficient bounds of \cref{thm:positive-current-order-zero} bound the variation of each coefficient of $S_l-S_k$ by a fixed multiple of its trace. Hence every coefficient is Cauchy in variation on compact sets and defines the positive limit $S$. Conversely a locally finite positive limit bounds the traces of all $S_k$. Every bounded product here annihilates pluripolar sets by \cref{prop:bt-poles-zero-capacity,prop:bt-capacity-basic}, so its increasing limit does too. Its restriction to $O_k$ is exactly $S_k$; it is zero on $U\setminus\bigcup_kO_k$, a pluripolar set. Finally $\mathbf1_{E_k}S_k=\mathbf1_{E_k}S$. Dominated convergence against the locally finite coefficient variations proves the last assertion. The trace determines their null sets.
\end{proof}

\begin{definition}\label{def:nonpluripolar-product}
When \eqref{eq:nonpluripolar-local-mass-condition} holds, the limit is the \emph{non-pluripolar product}
\[
 \langle\ddc u_1\wedge\cdots\wedge\ddc u_p\rangle=S.
\]
The angle brackets distinguish this construction from a product that retains polar mass. We omit them for bounded potentials, for which it is the Bedford--Taylor product. The empty product is $1$.
\end{definition}
\index{non-pluripolar product!truncation construction}

\begin{proposition}\label{prop:nonpluripolar-local-calculus}
Local finiteness and the product depend only on the currents $\ddc u_j$. The product is symmetric, local in the plurifine topology and multilinear in its positive current factors. For addition in one factor, the product of the sum is locally finite if and only if both summand products are locally finite. These assertions allow the local products of arbitrary closed positive $(1,1)$-currents to glue whenever their mass condition holds.
\end{proposition}
\begin{proof}
Symmetry holds at every bounded truncation. For locality, two tuples agreeing on a plurifine open $O$ have equal bounded products on
$O\cap\bigcap_j\{u_j>-k\}\cap\bigcap_j\{v_j>-k\}$.
These sets exhaust $O$ except for their pluripolar pole sets. Both limits give those sets zero mass, so the restrictions agree. The quasi-Lindel\"of argument in \cref{thm:bt-plurifine-locality} makes these restrictions measurable for the completed coefficient measures.

If $\ddc u_j=\ddc v_j$, their difference $w_j$ is smooth pluriharmonic: the trace equation and Weyl's lemma first make the distribution smooth, and its Hessian is zero. On a neighborhood $V$ of a compact test set choose $C$ with $|w_j|\leq C$. On $O_k(v)\cap V$, the bounded psh functions $\max\{v_j,-k\}$ and $\max\{u_j,-k-C\}+w_j$ coincide. Plurifine locality therefore identifies the restricted product with the corresponding restriction of the $u$-product at level $k+C$. Its trace is bounded by the mass of that product on a slightly larger compact neighborhood. Thus finiteness for $u$ implies finiteness for $v$; the reverse implication uses $-w_j$. Exhausting their common finite locus proves equality. Constants and real truncation levels cause no change, since a real level is bounded between adjacent integer levels.

\Needspace{6\baselineskip}
For additivity compare the positive bounded products on
\[
 E_k=\{u_1>-k\}\cap\{v_1>-k\}
                    \cap\bigcap_{j=2}^p\{u_j>-k\}.
\]
On this set $u_{1,k}+v_{1,k}=u_1+v_1$, whose truncation at $-2k$ has the same Hessian. Bounded multilinearity gives the sum of the two summand products there. For the other factors use level $2k$ on the sum side; they agree with their level-$k$ truncations on $E_k$. The sets $E_k$ increase and miss only the union of the pole sets. By \cref{prop:finite-locus-truncation}, each tuple has an increasing trace-measure limit even before local finiteness is known; this extended measure need not define a current. It annihilates every pluripolar set. Restricting each extended trace to $E_k$ gives its corresponding bounded-level trace, so the bounded identity just proved holds there. Monotone convergence as $k\to\infty$ gives equality of the extended traces on every Borel set. Positivity now shows that the uniform local mass condition for the sum is equivalent to the two separate conditions. Passing to the limits then gives equality of currents. Positive scaling follows by rescaling the levels, proving multilinearity. Ordinary open locality and potential independence give the gluing assertion. Local positive-current potentials exist by \cref{thm:local-positive-current-potential}.
\end{proof}

\begin{example}\label{ex:nonpluripolar-infinite-local-mass}
On a small bidisc write $L=-\log|y|$ and
\[
 u(x,y)=(|x|^2-1)\sqrt L\quad(y\ne0),\qquad u(x,0)=-\infty.
\]
Take $|x|^2<1/2$ and $|y|<\mathrm e^{-1}$. Off $y=0$ its complex Hessian has
\[
 u_{x\bar x}=\sqrt L,\qquad
 u_{y\bar y}=\frac{1-|x|^2}{16|y|^2L^{3/2}},\qquad
 u_{x\bar y}=-\frac{\bar x}{4\bar y\sqrt L},\qquad
 \det(u_{j\bar k})=\frac{1-2|x|^2}{16|y|^2L}.
\]
It is positive definite there. Bounded-above analytic removal in \cref{thm:psh-bounded-zero-extension} gives exactly the displayed value on the axis, so $u$ is psh. The smooth measure $(\ddc u)^2$ has a positive fixed multiple of this determinant as density. Its integral near the axis diverges, since normal radial integration contains $\int_0^r\mathrm dt/(t(-\log t))$. Truncation agrees with that smooth product off the axis, which the finite levels exhaust. Thus the local mass condition fails. Local non-pluripolar products of arbitrary potentials need not be finite.
\end{example}

\section{Closedness and cutoff gradients}
\label{sec:nonpluripolar-closedness}
The level indicators are not differentiable. We replace them by bounded psh cutoffs whose gradient energies are controlled by the restricted masses. Throughout this section $\mathrm d^{\mathrm c}=(\partial-\bar\partial)/(4\pi\mathrm i)$, as in \cref{sec:bt-monotone}; hence $\mathrm d\mathrm d^{\mathrm c}=\ddc$.

\Needspace{23\baselineskip}
\begin{lemma}\label{lem:bounded-psh-current-chain-rule}
Let $\Theta$ be a closed positive current of bidimension $(1,1)$ on a coordinate domain, and let $v$ be locally bounded psh. The current $\mathrm d(v\Theta)$ has measure coefficients, denoted $\mathrm dv\wedge\Theta$. There is a positive gradient measure
\begin{equation}\label{eq:bounded-current-gradient-measure}
 E_\Theta(v)=\mathrm dv\wedge\mathrm d^{\mathrm c}v\wedge\Theta
       =\tfrac12\ddc(v^2\Theta)-v\,\ddc v\wedge\Theta.
\end{equation}
Smooth decreasing approximations have gradients converging strongly in the coefficient-matrix energy seminorm on compact subsets. For every smooth real $f$, the bounded products satisfy
\begin{equation}\label{eq:bounded-current-chain-rule}
 \mathrm d(f(v)\Theta)=f'(v)\,\mathrm dv\wedge\Theta.
\end{equation}
For a compactly supported smooth real one-form $\eta$, a cutoff $\tau\geq0$ equal to one near its support, and a bounded Borel function $a$, there is a constant $C_\eta$ such that
\begin{equation}\label{eq:current-gradient-cauchy-schwarz}
 \left|\int a\eta\wedge\mathrm dv\wedge\Theta\right|^2
 \leq C_\eta\left(\int\tau E_\Theta(v)\right)
                 \left(\int_{\operatorname{supp}\eta}|a|^2\,
                                         \Theta\wedge\beta\right).
\end{equation}
\end{lemma}
\begin{proof}
All assertions are local near a compact test support. Add a constant to make $v$ and its smooth radial approximations $v_l\downarrow v$ nonnegative and uniformly bounded there. For a smooth real function $w$, the differential product rule gives
\[
 \mathrm dw\wedge\mathrm d^{\mathrm c}w\wedge\Theta
       =\tfrac12\ddc(w^2\Theta)-w\ddc w\wedge\Theta.
\]
The measure on the left is positive. Applied to $w=v_l-v_m$ and tested with $\tau$, this identity gives energy tending to zero as $l,m\to\infty$. Indeed \cref{cor:bt-polynomial-weights} makes the squared-weight term tend to zero after integration by parts against $\ddc\tau$, and each of the four expanded Hessian terms tends to the same integral with weight $v$ and Hessian $\ddc v$. Their signs cancel. For two independent indices the convergence is uniform on tails: otherwise select violating pairs with both indices strictly increasing and apply that same decreasing-continuity theorem.

Here is the meaning of the energy seminorm, including when $\Theta$ is singular. Its positive coefficient matrix has bounded Radon--Nikodym entries relative to $\mu=\Theta\wedge\beta$, by \cref{thm:positive-current-order-zero}. Write that measurable Hermitian matrix as $H$. For this bidimension, its quadratic form is the density of $\mathrm i\xi\wedge\bar\xi\wedge\Theta$ for a constant $(1,0)$-form $\xi$. Positivity for a countable dense family of $\xi$ makes $H$ positive semidefinite almost everywhere. In coordinates the gradient energy is the squared $L^2(\mu)$ norm of $H^{1/2}\partial w$, with the fixed normalization absorbed into $H$. The preceding estimate makes these vectors for $v_l$ Cauchy on compact subsets. Their strong limit defines the gradient modulo the zero seminorm. Multiplication by a bounded measurable scalar respects that quotient. Pointwise matrix Cauchy--Schwarz bounds the derivative paired with $a\eta$ by the two norms in \eqref{eq:current-gradient-cauchy-schwarz}. It also bounds every derivative coefficient in variation. Distributional convergence of $v_l\Theta$ to $v\Theta$, by dominated convergence, identifies that measure derivative with $\mathrm d(v\Theta)$.

The squared norms converge in $L^1(\mu)$, so their limit is a positive measure. Polynomial weighted continuity identifies its distribution with the right side of \eqref{eq:bounded-current-gradient-measure}. This also proves independence of the smooth approximations. Finally $f(v_l)\Theta\to f(v)\Theta$ by dominated convergence. The products $f'(v_l)$ times the gradients converge strongly: subtract the limiting gradient first, then use dominated convergence on its squared energy with $f'(v_l)\to f'(v)$. Passing through the smooth chain rule proves \eqref{eq:bounded-current-chain-rule}. A constant shift only translates the argument of $f$ and does not change the gradients, completing the proof.
\end{proof}

\begin{theorem}\label{thm:nonpluripolar-product-closed}
Every locally finite non-pluripolar product is closed.
\end{theorem}
\begin{proof}
Closedness is local. Work on a smaller coordinate neighborhood and subtract constants so that all local potentials $u_j\leq0$. If $p=n$, closedness holds by degree. Otherwise wedge the truncated products with a constant strongly positive $(n-p-1,n-p-1)$-form $\gamma$ and write
\[
 \Theta_k=\gamma\wedge\bigwedge_j\ddc u_{j,k},\qquad
 S=\gamma\wedge\langle\bigwedge_j\ddc u_j\rangle.
\]
These $\Theta_k$ are closed positive currents of bidimension $(1,1)$, and $\mathbf1_{O_k}\Theta_k\uparrow S$ with uniformly bounded restricted mass on compact sets.

Set $u=\sum_j u_j$. Choose a smooth convex increasing $\chi$ with $\chi=0$ on $(-\infty,1/2]$ and $\chi(1)=1$, and put $w_k=\chi(\mathrm e^{u/k})$. It is bounded psh, $0\leq w_k\leq1$, increases to one at every finite point of $u$, and is zero outside $O_k$: there $u\leq-k$ and $\mathrm e^{-1}<1/2$. At the poles define its value to be zero. Composition followed by the canonical representative gives precisely these pointwise values. Choose smooth $0\leq\zeta\leq1$, zero on $(-\infty,1/2]$ and equal to one near $1$. Both $\zeta(w_k)$ and $\zeta'(w_k)$ vanish outside $O_k$. Dominated convergence against $S$ shows
\[
 \zeta(w_k)\Theta_k\longrightarrow S,
 \qquad \zeta'(w_k)\longrightarrow0
                         \quad S\text{-almost everywhere}.
\]
For the first assertion use $\mathbf1_{O_k}\Theta_k=S$ on $O_k$ and that $S$ gives the poles zero mass.

For a compactly supported one-form $\eta$ and a cutoff $\tau\geq0$ equal to one near its support, the gradient identity and positivity give
\[
 \int\tau E_{\Theta_k}(w_k)
 \leq\tfrac12\int w_k^2\ddc\tau\wedge\Theta_k
 \leq C\int_{\operatorname{supp}\tau}
                           \mathbf1_{O_k}\Theta_k\wedge\beta\leq C'.
\]
The first inequality drops $w_k\ddc w_k\wedge\Theta_k\geq0$. The next uses $w_k=0$ off $O_k$, its bound by one, and the coefficient bound for the fixed smooth $\ddc\tau$. The mass hypothesis makes $C'$ independent of $k$. Apply \eqref{eq:current-gradient-cauchy-schwarz} with $a=\zeta'(w_k)$. Its second factor tends to zero, since
\[
 |\zeta'(w_k)|^2\Theta_k\wedge\beta
       =|\zeta'(w_k)|^2\mathbf1_{O_k}\Theta_k\wedge\beta
       \leq |\zeta'(w_k)|^2 S\wedge\beta,
\]
and the latter tends to zero by dominated convergence. The chain rule now gives $\mathrm d(\zeta(w_k)\Theta_k)\to0$, hence $\mathrm dS=0$. Constant strongly positive $\gamma$ span the real forms of this bidegree; wedging them with coordinate one-forms spans the test forms for the two types of $\mathrm d\langle\bigwedge_j\ddc u_j\rangle$. Thus the original product is closed.
\end{proof}
\index{non-pluripolar product!closedness}

\section{Global finiteness and bounded cohomological mass}
\label{sec:nonpluripolar-global-mass}
Let $(X,\omega)$ now be connected compact K\"ahler. For a closed smooth real $(1,1)$-form $\theta$ and a locally bounded $\theta$-psh potential $\varphi$, the local bounded products of $\theta+\ddc\varphi$ glue by smooth-potential independence. Smooth closed forms can be wedged with these currents.

\begin{proposition}\label{prop:bounded-global-cohomological-mass}
Let $\varphi_j\in\PSH(X,\theta_j)$ be locally bounded, $1\leq j\leq p\leq n$, and let $\eta$ be a smooth closed $(n-p,n-p)$-form. Then
\[
 \int_X\bigwedge_{j=1}^p(\theta_j+\ddc\varphi_j)\wedge\eta
       =\int_X\theta_1\wedge\cdots\wedge\theta_p\wedge\eta.
\]
\end{proposition}
\begin{proof}
For any closed positive remaining bounded product $R$, local recursion gives
$(\theta_j+\ddc\varphi_j)\wedge R=\theta_j\wedge R+\ddc(\varphi_jR)$.
Both sides glue, since $\varphi_j$ is a global bounded function. The derivative term integrates to zero against any smooth closed complementary form: its distributional pairing is with that form's zero Hessian. Replace the first factor, then the second, and so on, absorbing the already replaced smooth closed forms into the test form. All remaining nonsmooth products are positive and closed. The last integral is the displayed smooth one.
\end{proof}

\begin{theorem}\label{thm:compact-nonpluripolar-finiteness}
Every tuple of closed positive $(1,1)$-currents $T_1,\ldots,T_p$, $1\leq p\leq n$, on a compact K\"ahler manifold has a finite non-pluripolar product. Write $T_j=\theta_j+\ddc\varphi_j$ and choose $C\geq0$ so that $\tau_j=\theta_j+C\omega$ are positive definite. Then
\begin{equation}\label{eq:shifted-cohomological-mass-bound}
 \int_X\langle T_1\wedge\cdots\wedge T_p\rangle\wedge\omega^{n-p}
       \leq\int_X\tau_1\wedge\cdots\wedge\tau_p\wedge\omega^{n-p}.
\end{equation}
If the chosen $\theta_j$ are already semipositive, the same bound holds with $C=0$.
\end{theorem}
\begin{proof}
The global potential decomposition is \cref{thm:global-ddc-lemma}. Since $\tau_j$ is positive and $\varphi_j$ is $\tau_j$-psh, the global bounded potentials $\max\{\varphi_j,-k\}$ give positive currents
\[
 R_{j,k}=\tau_j+\ddc\max\{\varphi_j,-k\}.
\]
Their full mixed trace mass is the right side of \eqref{eq:shifted-cohomological-mass-bound}, independent of $k$, by \cref{prop:bounded-global-cohomological-mass}.

To bound the original local truncations, choose on a relatively compact coordinate chart smooth potentials $h_j$ for $\theta_j$ and $h$ for $\omega$. There $u_j=\varphi_j+h_j$ is psh. Choose $B$ with $h_j\leq B$ on a neighborhood of the compact test set. On $O_l(u)$, $\varphi_j>-l-B$. For any integer $k\geq l+B$, plurifine locality identifies $R_{j,k}$ there with the Hessian of $\max\{u_j,-l\}+Ch$: both local bounded psh potentials equal $u_j+Ch$ on that set. Their mixed bounded products agree there. Expanding the latter positive factors shows
\[
 \mathbf1_{O_l(u)}\bigwedge_j\ddc\max\{u_j,-l\}
       \leq\mathbf1_{O_l(u)}\bigwedge_j R_{j,k}.
\]
The global mass bound therefore implies every local mass condition. The products glue and are closed by the preceding results. The same argument for $T_j+C\omega$ identifies their non-pluripolar product with the increasing common-global-level restrictions of $\bigwedge_jR_{j,k}$; these levels exhaust all finite points of the $\varphi_j$. It has mass at most the same smooth bound. Multilinearity and positivity give
$\langle\bigwedge_j T_j\rangle\leq\langle\bigwedge_j(T_j+C\omega)\rangle$,
proving \eqref{eq:shifted-cohomological-mass-bound}. For semipositive $\theta_j$, the constant potentials and their maxima are $\theta_j$-psh, so the argument applies directly with $C=0$.
\end{proof}
\index{non-pluripolar product!global finiteness}

\begin{example}\label{ex:nonpluripolar-divisor-zero-mass}
On $\mathbb P^1$ the current $[0]$ has non-pluripolar product $\langle[0]\rangle=0$. Near zero its potential is $\log|z|^2$. The Hessian of its truncation is supported on the boundary level, excluded by $\{\log|z|^2>-k\}$. Away from zero the current is zero. Its ordinary mass is one, while its non-pluripolar mass is zero. Thus discarding poles can change even a first-degree cohomology class. In dimension zero the empty product is the counting measure and every positive-degree product is zero.
\end{example}

\section{Fibre integration and modifications}
\label{sec:nonpluripolar-fibres}
An integral on a product manifold can be disintegrated along its complex fibres. The bounded formula must hold for Borel level indicators before it can be passed through non-pluripolar truncation.

\begin{lemma}\label{lem:bounded-product-fibre-fubini}
Let $Q\subset\mathbb C^r$ and $D\subset\mathbb C^m$ be domains, with projection $p_D:Q\times D\to D$. Let $u_1,\ldots,u_r$ be locally bounded psh on $Q\times D$, and put
\[
 M=\ddc u_1\wedge\cdots\wedge\ddc u_r,\qquad
 M_z=\bigwedge_{j=1}^r\ddc_x u_j(\,\cdot\,,z).
\]
For a smooth real $(m,m)$-form $\eta$ on $D$ and a bounded Borel function $h$ of compact support, the fibre integral is Borel measurable and
\begin{equation}\label{eq:bounded-product-fibre-fubini}
 \int_{Q\times D}h\,p_D^*\eta\wedge M
       =\int_D\left(\int_Q h(x,z)\,M_z(x)\right)\eta(z).
\end{equation}
For nonnegative $h$ and $\eta$ the same identity holds with possibly infinite integrals.
\end{lemma}
\begin{proof}
First take $\eta\geq0$ and $h$ continuous of compact support. Smooth decreasing regularizations of the $u_j$ exist on a product neighborhood of that support. In the smooth exterior product, wedging with $p_D^*\eta$ kills every term containing a base differential. Ordinary Fubini therefore gives \eqref{eq:bounded-product-fibre-fubini} for the regularizations. Decreasing Bedford--Taylor continuity gives weak convergence of the products on the total space and on every fibre. The common bounds on the potentials and \cref{thm:cln-bounded} give a uniform fibre mass bound on the compact projection of the support. Dominated convergence proves the continuous-test formula. A finite cover and a partition of unity handle supports requiring more than one such neighborhood.

To justify the Borel tests, fix $U\Subset Q$ and $V\Subset D$. For $a\in C_c(U)$, $z\mapsto\int a\,M_z$ is Borel: it is the pointwise limit of the corresponding continuous smooth-regularization integrals. Every relatively open subset $A\subset U$ admits continuous functions $0\leq a_l\leq1$ of compact support in $U$ increasing to $\mathbf1_A$; construct them from distance to the complement and a compact exhaustion. Monotone convergence proves measurability of $M_z(A)$, including $A=U$. All these measures are finite on $U$, uniformly for $z\in V$. The sets $A$ for which $M_z(A)$ is Borel as a function of $z$ are therefore closed under complements in $U$ and countable disjoint unions. The intersection-closed open sets generate the Borel sets, so the class argument proved in \cref{prop:tonelli-fubini} gives measurability for every Borel $A$.

Apply the same argument to sets $E\subset U\times V$: the class for which $z\mapsto M_z(E_z)$ is Borel contains the Borel rectangles and is closed under complements and countable disjoint unions. It contains all Borel sets. Increasing exhaustions of $Q$ and $D$ give the assertion on the whole product. Hence
\[
 \lambda(E)=\int_D M_z(E_z)\,\eta(z)
\]
is a Borel measure. Countable additivity follows from monotone convergence on each fibre and then on $D$. The uniform local fibre mass bound makes it locally finite, and \cref{prop:borel-regularity} makes it Radon. It agrees with $p_D^*\eta\wedge M$ on continuous compactly supported tests by the first paragraph. Radon uniqueness makes the measures equal. Simple approximation proves the formula for nonnegative Borel functions and then for bounded signed compactly supported functions. On a compact base neighborhood an arbitrary smooth real top form is a difference of two smooth nonnegative top forms, which proves the signed assertion. The cases $r=0$ or $m=0$ reduce to counting measure and the empty product.
\end{proof}
\index{fibre integration!Bedford--Taylor products}

\begin{corollary}\label{cor:compact-nonpluripolar-fibre-fubini}
Let $X,Z$ be connected compact K\"ahler manifolds of dimensions $r,m$, with projections $p_X,p_Z$. Let $\theta$ be a closed smooth real $(1,1)$-form on $X$, $\varphi\in\PSH(X,\theta)$, and let $\gamma\geq0$ be a closed smooth $(1,1)$-form on $Z$. Set
\[
 \Theta=p_X^*\theta+p_Z^*\gamma,\qquad
 \Phi=p_X^*\varphi,\qquad
 T=\Theta+\ddc W,\quad S=\Theta+\ddc\Phi,
\]
where $W\in\PSH(X\times Z,\Theta)$ and $\Phi-C\leq W\leq\Phi+C$ for a finite $C$. Every $W_z=W(\,\cdot\,,z)$ belongs to $\PSH(X,\theta)$. For $0\leq j\leq r$ and any smooth real $(m,m)$-form $\eta$ on $Z$,
\begin{equation}\label{eq:compact-nonpluripolar-fibre-fubini}
 \begin{aligned}
 &\int_{X\times Z}(W-\Phi)\,p_Z^*\eta\wedge
                              \langle T^j\wedge S^{r-j}\rangle\\
 &\qquad=\int_Z\left(\int_X(W_z-\varphi)
            \langle(\theta+\ddc W_z)^j\wedge
                         (\theta+\ddc\varphi)^{r-j}\rangle\right)\eta(z).
 \end{aligned}
\end{equation}
The same identity holds with both difference weights replaced by $1$. The difference weights may be assigned any bounded value at their common poles.
\end{corollary}
\Needspace{6\baselineskip}
\begin{proof}
In product charts choose smooth local potentials $a(x)$ for $\theta$ and $b(z)$ for $\gamma$. The restrictions of $W+a+b$ to fibres are psh, and $W_z$ is not identically $-\infty$ because $W_z\geq\varphi-C$ and $\varphi$ is finite somewhere. Thus the asserted fibre potentials are defined. All products are finite by \cref{thm:compact-nonpluripolar-finiteness}. That theorem also bounds their fibre masses uniformly in $z$: the shifted smooth background can be chosen from the fixed $\theta$, independently of $W_z$.

Put $w=W+a+b$, $v=\Phi+a+b$ and $O_k=\{w>-k\}\cap\{v>-k\}$. Apply \eqref{eq:bounded-product-fibre-fubini} to the bounded truncations of $w,v$, with multiplicities $j,r-j$, and the Borel weight $\rho G\mathbf1_{O_k}$, where $\rho\geq0$ is a smooth compactly supported chart cutoff and $G\geq0$ is bounded Borel. Their restricted mixed measures increase to the non-pluripolar products on the total space and on every fibre. When one multiplicity is zero, the additional finite-level condition still exhausts that product: its omitted pole set is pluripolar and has zero mass. Monotone convergence gives the same fibre formula with the non-pluripolar limits and weight $\rho G$. Take $G=1$ for the unweighted formula and $G=W-\Phi+C$ on their common finite locus for the weighted formula; assign $G=C$ at their poles. Subtract the constant-weight identity. Sum over a finite nonnegative partition of unity and use linearity for signed $\eta$. The uniform mass bound and $|W-\Phi|\leq C$ justify every finite signed integral. Polar sets have zero product mass, so the assigned values do not affect them.
\end{proof}
\index{fibre integration!non-pluripolar products}

\begin{proposition}\label{prop:nonpluripolar-modification-pushforward}
Let $\pi:\widetilde X\to X$ be a modification between compact K\"ahler manifolds and let $T_1,\ldots,T_p$ be closed positive $(1,1)$-currents on $X$, $1\leq p\leq\dim X$. Then
\[
 \pi_*\langle\pi^*T_1\wedge\cdots\wedge\pi^*T_p\rangle
                  =\langle T_1\wedge\cdots\wedge T_p\rangle.
\]
In particular their top-degree masses agree.
\end{proposition}
\begin{proof}
Pullback is defined by pulling back local psh potentials; their pluriharmonic differences pull back as well. A potential cannot become identically $-\infty$ on a component, because a modification is a biholomorphism over a dense open set. Both products are finite. Choose a proper analytic set $A\subset X$ over whose complement $\pi$ is a biholomorphism, as in the definition of modification. Potential independence and coordinate invariance identify the products over $X\setminus A$. The downstairs product gives $A$ zero mass. The upstairs product gives the proper analytic set $\pi^{-1}(A)$ zero mass, so its pushforward also gives $A$ zero coefficient variation. Equality off $A$ therefore extends to all coefficient measures. Testing the constant function in top degree proves the last assertion.
\end{proof}
\index{non-pluripolar product!modification pushforward}
